Recordemos que $A^*$ esta dado por
\[f(n)=g(n)+h(n)\]
Y además una heurística admisible es aquella la cual se cumple que:
\[h(n) \leq h^*(n)\]
Es decir, si el costo estimado es menor o igual al costo real mínimo.\\
\textbf{Demostración}\\
Sea $h$ una heurística admisible.  Probaremos por contradicción, suponiendo que $A^*$ no encontrará una solución óptima, entonces tenemos un camino subóptimo $C'$ que llega a la meta $s_f$ (por lo que $h(s_f)=0$) donde $g(s_f) > g(C^*)$ siendo $C^*$ un camino óptimo, entonces tenemos que

\[ f(s_f)= g(s_f) + h(s_f) = g(s_f) \]

Por lo que tenemos $f(s_f) > g (C^*)$, veamos que $A^*$ no completó a $C^*$ por lo que existe un estado $s$ que pertenece a dicho camino el cual fue abierto pero no expandido, con lo que se tiene que al ser parte del camino óptimo sabemos que

\[ g(s)+h^*(s) = g(C^*) \]
Donde $h^*(s)$ es el costo real mínimo desde $s$ hasta $s_f$.

Veamos entonces que al sumar $g(s)$ a la desigualdad dada por la admisibilidad tenemos que
\begin{align*}
    h(s) & \leq h^*(s)\\
    g(s) + h(s) & \leq g(s) + h^*(s)\\
    f(s) & \leq g(C^*)
\end{align*}

Entonces notemos que 
\begin{align*}
    g(C^*) & < f(s_f)\\
    f(s) & \leq g(C^*) 
\end{align*}

Y por tanto
\[f(s) \leq g(C^*) < f(s_f)\]
\[f(s) < f(s_f) \]
Esto implica que $A^*$ debió expandir a $s$ por lo que ocurre una contradicción.\\
$\therefore$ Si $h(s_t),\, s_t \in S$, es una heurística admisible para el algoritmo de $A^*$, entonces el algoritmo de $A^*$ encontrará una solución óptima. \textbf{QED}. 
 
